\documentclass[11pt]{article}

\usepackage[a4paper,margin=1in]{geometry}
\usepackage{amsmath,amssymb,bm}
\usepackage{booktabs}
\usepackage{array}
\usepackage{graphicx}
\usepackage{hyperref}
\usepackage{float}

\title{Bayesian Inference of Stochastic Gravitational-Wave Anisotropies with BLIP}
\author{Author Name(s)}
\date{\today}

\begin{document}
\maketitle

\begin{abstract}
This paper documents a Bayesian anisotropy-analysis pipeline for the Laser Interferometer Space Antenna (LISA) as implemented in the Bayesian LISA Inference Package (BLIP). The method models the measured time-delay interferometry (TDI) cross-spectra as the sum of instrumental noise and an anisotropic stochastic gravitational-wave background (SGWB), with the signal covariance built from the LISA response to a sky distribution expanded in spherical harmonics. The implementation supports both a full spherical-harmonic anisotropy model, in which sky structure is inferred through harmonic coefficients, and a reduced single-multipole total-power model, in which only the total power in one fixed multipole $L$ is sampled. We present the detector model, the spectral and spatial parameterizations, the prior transforms, the covariance likelihood, and the sampler choices used by BLIP. We also provide a results section with explicit placeholders for posterior figures, evidence tables, and sky-map summaries that can be filled once production runs are complete.
\end{abstract}

\section{Introduction}
The angular structure of a stochastic gravitational-wave background carries astrophysical and cosmological information beyond what is available from an isotropic monopole alone. For LISA, anisotropies can arise from unresolved Galactic binaries, nearby satellite populations, or genuinely anisotropic cosmological backgrounds. Recovering this information requires a forward model that connects sky power to the time-dependent response of the moving detector and a statistical framework that propagates instrumental uncertainties, spectral uncertainties, and sky uncertainties simultaneously.

BLIP is designed precisely for this task. It combines a generative model for LISA TDI data, configurable SGWB spectral models, anisotropic detector response calculations, and Bayesian posterior sampling. In the code base used here, anisotropy inference is available in two related forms. The first is the general spherical-harmonic mode, which samples coefficients describing anisotropic structure and reconstructs a sky map. The second is a reduced single-multipole mode, intended as a first-pass search, which fixes one multipole $L$ and infers only its total power amplitude and spectral slope. The latter does not reconstruct individual $b_{\ell m}$ coefficients and does not produce a sky map.

This manuscript is written to match the current BLIP implementation. The goal is not only to describe the conceptual method, but also to record what the code actually does physically and mathematically: how the time-domain injections are generated, how the frequency-domain data products are built, how the anisotropic response basis is computed, how priors are applied from configuration files, how the likelihood is assembled from covariance matrices, and how the posterior is sampled. The presentation below is therefore suitable both as a methods section for a paper and as technical documentation for reproducibility.

\section{Methods}
\subsection{Analysis Configuration and Data Products}
BLIP is configured from an \texttt{ini} file specifying the observation duration $T_{\mathrm{obs}}$, segment length $T_{\mathrm{seg}}$, sampling rate $f_s$, analysis band $[f_{\min},f_{\max}]$, TDI basis, LISA configuration, prior bounds, and sampler settings. In a representative fixed-multipole production run, the code uses
\begin{equation}
f_{\min}=4\times10^{-4}\ \mathrm{Hz}, \qquad
f_{\max}=3\times10^{-3}\ \mathrm{Hz}, \qquad
T_{\mathrm{obs}}=2\times10^{6}\ \mathrm{s}, \qquad
T_{\mathrm{seg}}=10^{5}\ \mathrm{s},
\end{equation}
with the \texttt{xyz} TDI channels, \texttt{nside}=4 for the sky integration, and a fixed multipole choice such as $L=2$.

The code can either load external data or simulate it. When external time series are loaded, BLIP optionally converts Doppler observables to strain-like variables. When injections are simulated, the code constructs three TDI channels, adds instrumental noise, and then adds SGWB realizations consistent with the model covariance. The characteristic LISA transfer-frequency scale is
\begin{equation}
f_{\star}=\frac{c}{2\pi L_{\mathrm{arm}}},
\qquad
f_0=\frac{f}{2f_{\star}}=\frac{\pi L_{\mathrm{arm}} f}{c},
\end{equation}
with $L_{\mathrm{arm}}=2.5\times10^9\ \mathrm{m}$ in the current implementation.

The likelihood is evaluated on Fourier-domain data products. For each channel $I\in\{1,2,3\}$ the code computes Fourier coefficients $r_I(f_k,t_n)$ from windowed time segments. A Hann window is applied and the Fourier coefficients are normalized so that $|r_I|^2$ has the units of a one-sided power spectral density. The data matrix used in the likelihood is the segment-wise cross-spectrum
\begin{equation}
D_{IJ}(f_k,t_n)=r_I^{\ast}(f_k,t_n)\,r_J(f_k,t_n).
\end{equation}
In the current implementation the number of likelihood segments is
\begin{equation}
N_{\mathrm{seg}}=\left\lfloor \frac{T_{\mathrm{obs}}}{T_{\mathrm{seg}}}\right\rfloor-1,
\end{equation}
and the segment loop advances by one full segment at a time. This is worth stating explicitly because comments in the source mention $50\%$ overlap, whereas the implemented FFT loop currently uses adjacent, non-overlapping likelihood segments.

\subsection{Instrumental Noise Model}
BLIP models the fundamental LISA noise budget using one position-noise scale $N_p$ and one acceleration-noise scale $N_a$. The corresponding one-link spectra are
\begin{align}
S_p(f) &= N_p\left[1+\left(\frac{2\times10^{-3}}{f}\right)^4\right], \\
S_a(f) &= N_a\left[1+\frac{16\times10^{-8}}{f^2}\right]
\left[1+\left(\frac{f}{8\times10^{-3}}\right)^4\right]
\left(\frac{1}{2\pi f}\right)^4 .
\end{align}
For Michelson variables the auto- and cross-spectra are implemented as
\begin{align}
S_{\mathrm{auto}}(f) &= 4\left[2S_a(f)\left(1+\cos^2 2f_0\right)+S_p(f)\right], \\
S_{\mathrm{cross}}(f) &= \left[-2S_p(f)-8S_a(f)\right]\cos(2f_0),
\end{align}
from which the full $3\times3$ Michelson covariance matrix is built. The \texttt{xyz} noise model is obtained by multiplying the Michelson response by $4\sin^2(2f_0)$, and the \texttt{aet} covariance is then formed by the standard linear combinations of the $X,Y,Z$ channels.

In terms of sampled parameters, BLIP works with $\log_{10}N_p$ and $\log_{10}N_a$. For any proposed point in parameter space, the instrumental covariance contribution is
\begin{equation}
\mathbf{C}_{\mathrm{noise}}(f_k,t_n)=\mathbf{C}_{\mathrm{noise}}(f_k),
\end{equation}
that is, stationary in time and repeated over all segments.

\subsection{Spectral Model for the SGWB}
The basic SGWB spectrum used in the anisotropy analyses considered here is a power law in the energy density,
\begin{equation}
\Omega(f)=A\left(\frac{f}{f_{\mathrm{ref}}}\right)^{\alpha},
\end{equation}
where $A$ is either the isotropic amplitude $\Omega_0$ or, in the single-multipole mode, the total-power amplitude $A_L$. The code also supports broken and truncated power laws, but the power-law case is the most relevant for the anisotropy examples presently configured.

BLIP converts $\Omega(f)$ to a strain power spectral density using
\begin{equation}
S_{\mathrm{gw}}(f)=\Omega(f)\,\frac{3}{4f^3}\left(\frac{H_0}{\pi}\right)^2,
\end{equation}
with $H_0=2.2\times10^{-18}\ \mathrm{s}^{-1}$ in the implementation. The total SGWB covariance is then the product of this scalar spectrum and a channel-space response matrix,
\begin{equation}
\mathbf{C}_{\mathrm{gw}}(f_k,t_n)=S_{\mathrm{gw}}(f_k)\,\mathbf{R}(f_k,t_n).
\end{equation}

\subsection{LISA Orbital Response and Anisotropy Basis}
The code evaluates the LISA constellation at the midpoint of each likelihood segment using the analytic MLDC orbit prescription. For \texttt{lisa\_config=orbiting}, the time dependence of the three spacecraft positions is retained; for \texttt{lisa\_config=stationary}, the geometry is frozen to the initial epoch. The anisotropic response is computed numerically by integrating over a HEALPix sky grid of resolution set by \texttt{nside}.

For each sky direction $\hat{\Omega}$, each time segment $t_n$, and each frequency bin $f_k$, BLIP evaluates the polarization-summed products of the LISA antenna pattern functions. The spherical-harmonic response basis is then constructed as
\begin{equation}
R_{IJ}^{LM}(f_k,t_n)
=\frac{\Delta\Omega}{2}\sum_{p\in\{+,\times\}}
\sum_{\hat{\Omega}}
F_I^{p}(\hat{\Omega},f_k,t_n)\,
F_J^{p\ast}(\hat{\Omega},f_k,t_n)\,
Y_{LM}(\hat{\Omega}),
\end{equation}
where $\Delta\Omega$ is the HEALPix pixel area and the sum over sky direction is the discrete approximation to the full sky integral. The resulting object is a complex Hermitian $3\times3$ covariance basis for each $(f_k,t_n)$ and each harmonic index $(L,M)$.

\subsection{Full Spherical-Harmonic Anisotropy Model}
In the full anisotropic mode, BLIP does not sample the power-map coefficients $a_{LM}$ directly. Instead it samples a set of coefficients $b_{\ell m}$ and converts them to $a_{LM}$ through Clebsch--Gordan couplings,
\begin{equation}
a_{LM}=\sum_{\ell_1 m_1}\sum_{\ell_2 m_2}
\beta^{LM}_{\ell_1 m_1,\ell_2 m_2}\,
b_{\ell_1 m_1}\,b_{\ell_2 m_2},
\end{equation}
where the coupling coefficients $\beta^{LM}_{\ell_1 m_1,\ell_2 m_2}$ are precomputed from Clebsch--Gordan coefficients. Because the power map is quadratic in the sampled $b_{\ell m}$ coefficients, the maximum harmonic order of the power map is $L_{\max}=2\ell_{\max}$ in this mode.

After constructing the $a_{LM}$, the code normalizes them according to
\begin{equation}
\tilde{a}_{LM}=\frac{a_{LM}}{a_{00}\sqrt{4\pi}},
\end{equation}
so that the monopole is fixed to the conventional normalization. The sky-convolved response is then
\begin{equation}
R_{IJ}(f_k,t_n\mid \bm{b})
=\sum_{LM} R_{IJ}^{LM}(f_k,t_n)\,\tilde{a}_{LM}.
\end{equation}
Finally, the anisotropic SGWB covariance becomes
\begin{equation}
C_{IJ}^{\mathrm{asgwb}}(f_k,t_n\mid \alpha,A,\bm{b})
=S_{\mathrm{gw}}(f_k\mid \alpha,A)\,
R_{IJ}(f_k,t_n\mid \bm{b}).
\end{equation}

\subsection{Single-Multipole Total-Power Model}
The simplified anisotropy mode implemented for first-pass searches keeps one chosen multipole $L$ fixed and marginalizes over its $m$-modes analytically rather than sampling individual coefficients. In this mode the sampled parameters are just
\begin{equation}
\left\{\log_{10}N_p,\ \log_{10}N_a,\ \alpha,\ \log_{10}A_L\right\}.
\end{equation}
No sky map is reconstructed.

The code first evaluates the spherical-harmonic response basis up to the selected multipole and then extracts the subset corresponding to that single $L$. It combines the $m$-mode responses into an effective covariance response through the Gram-matrix sum
\begin{equation}
\mathbf{R}^{(L)}_{\mathrm{eff}}(f_k,t_n)
=\frac{1}{2L+1}\sum_{m=-L}^{L}
\mathbf{R}^{(Lm)}(f_k,t_n)\,
\mathbf{R}^{(Lm)\dagger}(f_k,t_n),
\end{equation}
which is the matrix form of the operation implemented in the code. This produces an explicitly Hermitian, positive semi-definite response matrix describing the total statistically isotropic power in that fixed multipole. The corresponding SGWB covariance is
\begin{equation}
\mathbf{C}_{\mathrm{gw}}^{(L)}(f_k,t_n)
=S_{\mathrm{gw}}(f_k\mid \alpha,A_L)\,
\mathbf{R}_{\mathrm{eff}}^{(L)}(f_k,t_n).
\end{equation}
This mode is therefore especially useful for detectability studies or for establishing whether a given angular scale can be constrained before undertaking a full sky reconstruction.

\subsection{Priors}
BLIP samples on the unit hypercube and applies deterministic prior transforms to map each cube coordinate $u\in[0,1]$ to a physical parameter:
\begin{equation}
x=x_{\min}+u\,(x_{\max}-x_{\min}).
\end{equation}
All bounds are configurable in the \texttt{[priors]} section of the \texttt{ini} file, but the implementation includes legacy defaults. Table~\ref{tab:priors} summarizes the main priors relevant for anisotropy analyses.

\begin{table}[H]
\centering
\caption{Uniform prior intervals used by BLIP for the anisotropy analyses. The rightmost column shows one representative tuned configuration for the single-multipole production run.}
\label{tab:priors}
\begin{tabular}{>{\raggedright\arraybackslash}p{0.22\linewidth} >{\raggedright\arraybackslash}p{0.33\linewidth} >{\raggedright\arraybackslash}p{0.18\linewidth} >{\raggedright\arraybackslash}p{0.18\linewidth}}
\toprule
Parameter & Meaning & Default interval & Example tuned interval \\
\midrule
$\log_{10}N_p$ & Position-noise level & $[-44,-39]$ & $[-43,-40]$ \\
$\log_{10}N_a$ & Acceleration-noise level & $[-51,-46]$ & $[-50,-47]$ \\
$\alpha$ & Power-law slope & $[-5,5]$ & $[-1,2.5]$ \\
$\log_{10}\Omega_0$ & Isotropic SGWB amplitude & $[-14,8]$ & placeholder \\
$\log_{10}A_L$ & Fixed-$L$ total-power amplitude & $[-14,8]$ & $[-9.5,-5.5]$ \\
$\alpha_1$ & Broken-power-law low-frequency slope & $[-4,6]$ & placeholder \\
$\alpha_2$ & Broken-power-law high-frequency slope & $[0,40]$ & placeholder \\
$\log_{10}f_{\mathrm{break}}$ & Break frequency & $[-4,-2]$ & placeholder \\
$\log_{10}f_{\mathrm{cut}}$ & Cutoff frequency & $[-4,-2]$ & placeholder \\
$\log_{10}f_{\mathrm{scale}}$ & Truncation scale & $[-4,-2]$ & placeholder \\
\bottomrule
\end{tabular}
\end{table}

For the full spherical-harmonic anisotropy model, the additional priors are imposed on the $b_{\ell m}$ parameterization. The coefficient $b_{00}$ is fixed to unity. For $m=0$, BLIP samples
\begin{equation}
b_{\ell 0}\sim\mathcal{U}(-3,3),
\end{equation}
while for $m>0$ it samples amplitude and phase separately as
\begin{equation}
|b_{\ell m}|\sim\mathcal{U}(0,3),
\qquad
\phi_{\ell m}\sim\mathcal{U}(-\pi,\pi),
\qquad
b_{\ell m}=|b_{\ell m}|\,e^{i\phi_{\ell m}}.
\end{equation}
These priors encode a broad but finite anisotropy space without imposing a preferred phase orientation.

\subsection{Covariance Likelihood}
For any proposed parameter vector $\bm{\theta}$, BLIP computes the total covariance as the sum of the contributions from each active submodel,
\begin{equation}
\mathbf{C}(f_k,t_n\mid\bm{\theta})
=\sum_a \mathbf{C}_a(f_k,t_n\mid\bm{\theta}_a).
\end{equation}
In the present use case this is typically the sum of an instrumental-noise covariance and one anisotropic SGWB covariance term.

The likelihood is the complex Gaussian likelihood for the Fourier-domain TDI data. Writing the data outer product as $\mathbf{D}(f_k,t_n)$, the code evaluates
\begin{equation}
\ln \mathcal{L}(\bm{\theta})
=-\sum_{k,n}\mathrm{Tr}\!\left[
\mathbf{C}^{-1}(f_k,t_n\mid\bm{\theta})\,
\mathbf{D}(f_k,t_n)
\right]
-\sum_{k,n}\ln\!\left[
\pi T_{\mathrm{seg}}\,
\left|\det \mathbf{C}(f_k,t_n\mid\bm{\theta})\right|
\right].
\label{eq:likelihood}
\end{equation}
This is exactly the structure implemented in BLIP. The matrix inverse and determinant are computed with a custom vectorized $3\times3$ routine rather than repeated generic linear-algebra calls, which improves performance because every likelihood evaluation requires inverses across all frequencies and time segments.

\subsection{Injection Generation}
When synthetic data are generated, the SGWB contribution is produced directly from the model covariance. For each splice segment the code performs a Cholesky decomposition of the channel covariance,
\begin{equation}
\mathbf{C}(f_k,t_n)=\mathbf{L}(f_k,t_n)\mathbf{L}^{\dagger}(f_k,t_n),
\end{equation}
draws complex standard-normal random vectors $\bm{z}$, and rescales them as
\begin{equation}
\tilde{\bm{h}}(f_k,t_n)=\mathbf{L}(f_k,t_n)\bm{z}.
\end{equation}
The resulting frequency-domain realization is inverse-Fourier transformed into time-domain splices, multiplied by a sine splice window, and overlap-added. In the current implementation the SGWB injection uses half-overlapping splices of fixed duration $10^4\ \mathrm{s}$. This injection-side overlap is separate from the likelihood-side Fourier segmentation described above.

\subsection{Samplers and Evidence Calculation}
BLIP currently supports three posterior engines.

\paragraph{Dynesty nested sampling.}
This is the default choice in the shipped anisotropy configuration files. The code instantiates \texttt{dynesty.NestedSampler} with \texttt{bound='multi'} and a user-selected proposal method such as \texttt{rwalk} or \texttt{rslice}. Posterior samples are obtained by equal-weight resampling of the dead-point archive, and the run returns both $\log Z$ and its estimated error.

\paragraph{Nessai nested sampling.}
BLIP also supports \texttt{nessai}, which performs nested sampling with normalizing-flow proposals. The code wraps the BLIP unit-cube prior and transformed likelihood in a \texttt{nessai.Model} adapter, uses a unit-cube bound for each parameter, and configures the flow width through the number of neurons. This backend likewise returns posterior samples and a Bayesian evidence estimate.

\paragraph{Emcee ensemble MCMC.}
For posterior exploration without evidence estimation, BLIP provides an \texttt{emcee} backend. It uses a stretch move on walkers initialized uniformly in the unit cube, checks that walkers remain in the cube, applies the BLIP prior transform, and evaluates the same covariance likelihood as Eq.~(\ref{eq:likelihood}). This backend yields posterior samples but not a nested-sampling evidence.

For reproducibility, BLIP can fix the random seed and checkpoint nested-sampling runs. Posterior samples are written to \texttt{post\_samples.txt}; nested-sampling evidences are stored in \texttt{logz.txt} and \texttt{logzerr.txt} where applicable.

\subsection{Outputs Used in the Results Section}
The default BLIP workflow writes the following products that are directly useful for a paper:
\begin{itemize}
\item posterior samples in \texttt{post\_samples.txt};
\item a serialized model in \texttt{model.pickle};
\item an injection object in \texttt{injection.pickle} for simulated runs;
\item posterior corner plots in \texttt{corners.png};
\item fit-diagnostic plots comparing the recovered covariance spectra to the injected or measured spectra;
\item sky maps when the model has a true anisotropic map degree of freedom.
\end{itemize}
These files provide the numerical values and figures that should replace the explicit placeholders in the next section.

\section{Results}
\subsection{Analysis Setup}
This subsection should summarize the exact configuration used for the production analysis.

\begin{table}[H]
\centering
\caption{Placeholder table for the final run configuration. Replace the entries below with the values used in the production analysis.}
\label{tab:runconfig}
\begin{tabular}{lll}
\toprule
Quantity & Symbol / setting & Value used in final run \\
\midrule
Observation duration & $T_{\mathrm{obs}}$ & placeholder \\
Segment length & $T_{\mathrm{seg}}$ & placeholder \\
Frequency band & $[f_{\min},f_{\max}]$ & placeholder \\
TDI basis & channels & placeholder \\
Sky resolution & \texttt{nside} & placeholder \\
Anisotropy model & full sph or fixed $L$ & placeholder \\
Multipole choice & $L$ (if applicable) & placeholder \\
Sampler & dynesty / nessai / emcee & placeholder \\
Sampler setting & $n_{\mathrm{live}}$ or walkers & placeholder \\
Seed & random seed & placeholder \\
\bottomrule
\end{tabular}
\end{table}

Placeholder text: describe whether the run was an injection-recovery study or an analysis of externally supplied data, identify the spectral model, state the prior ranges actually used, and mention whether the analysis was carried out in the full spherical-harmonic mode or in the single-multipole total-power mode.

\subsection{Posterior Constraints}
This subsection should report the posterior constraints on the sampled parameters. For the single-multipole mode, the key results are the posteriors on $\log_{10}N_p$, $\log_{10}N_a$, $\alpha$, and $\log_{10}A_L$. For the full anisotropic mode, it should also summarize the posterior support on the $b_{\ell m}$ coefficients or derived $a_{LM}$ values.

\begin{figure}[H]
\centering
\fbox{\parbox[c][2.2in][c]{0.82\linewidth}{\centering
Placeholder for the BLIP corner plot.\\
Insert \texttt{corners.png} or an updated publication-quality posterior figure here.
}}
\caption{Placeholder posterior summary figure. In the single-multipole analysis this figure should show the joint posterior for the noise parameters, spectral slope, and total multipole power. In the full spherical-harmonic analysis it should be supplemented by a compact summary of the spatial posterior.}
\label{fig:corner}
\end{figure}

\begin{table}[H]
\centering
\caption{Placeholder table for posterior summaries. Replace with medians, credible intervals, and injected values where relevant.}
\label{tab:posteriors}
\begin{tabular}{llll}
\toprule
Parameter & Prior range & Posterior summary & Truth / reference \\
\midrule
$\log_{10}N_p$ & placeholder & placeholder & placeholder \\
$\log_{10}N_a$ & placeholder & placeholder & placeholder \\
$\alpha$ & placeholder & placeholder & placeholder \\
$\log_{10}A_L$ or $\log_{10}\Omega_0$ & placeholder & placeholder & placeholder \\
Spatial parameter(s) & placeholder & placeholder & placeholder \\
\bottomrule
\end{tabular}
\end{table}

Placeholder text: discuss whether the posterior is prior-dominated or likelihood-dominated, whether the recovered amplitude excludes zero power in the chosen multipole, and whether the spectral slope is consistent with the injected or expected astrophysical value.

\subsection{Goodness of Fit in the TDI Spectra}
The next result to report is how well the recovered covariance reproduces the measured auto- and cross-channel spectra. This is the natural place to compare the recovered model against the injected spectra in simulation studies or against the measured spectra in real-data studies.

\begin{figure}[H]
\centering
\fbox{\parbox[c][2.2in][c]{0.82\linewidth}{\centering
Placeholder for PSD and covariance-fit diagnostics.\\
Insert the BLIP fitmaker outputs, data-model spectral overlays, or residual plots here.
}}
\caption{Placeholder fit diagnostic. This figure should compare the posterior-predicted TDI spectra against the data, ideally showing both auto-spectra and representative cross-spectra.}
\label{fig:fit}
\end{figure}

Placeholder text: quantify how well the posterior median model tracks the measured spectra across the analysis band, note any frequency regions where the fit degrades, and discuss degeneracies between noise parameters and anisotropic power.

\subsection{Anisotropy Interpretation}
The interpretation subsection depends on which anisotropy mode was used.

If the full spherical-harmonic mode was used, this subsection should present the recovered sky map or the posterior on the dominant $a_{LM}$ coefficients. It should state which multipoles are significantly constrained, whether the monopole-normalized sky map shows coherent angular structure, and how the inferred morphology compares to the injected or expected sky distribution.

If the single-multipole total-power mode was used, this subsection should instead interpret the inferred $A_L$ as a statement about the strength of anisotropy at angular scale $L$. In that case the recommended discussion points are:
\begin{itemize}
\item whether the posterior on $A_L$ is clearly separated from the lower prior boundary;
\item how the inferred power changes when the prior is widened or narrowed;
\item whether the data favor one multipole over another in model-comparison tests;
\item whether the detected angular power is robust to changes in the frequency band or noise priors.
\end{itemize}

\begin{figure}[H]
\centering
\fbox{\parbox[c][2.2in][c]{0.82\linewidth}{\centering
Placeholder for anisotropy-specific figure.\\
Use a sky map for the full spherical-harmonic mode, or a bar chart / posterior comparison across multipoles for the fixed-$L$ mode.
}}
\caption{Placeholder anisotropy summary figure.}
\label{fig:anisotropy}
\end{figure}

\subsection{Bayesian Evidence and Model Comparison}
Where nested sampling is used, BLIP provides the Bayesian evidence. This subsection should compare the evidence for competing models such as
\begin{equation}
\mathcal{M}_{\mathrm{noise}}, \qquad
\mathcal{M}_{\mathrm{noise+isotropic}}, \qquad
\mathcal{M}_{\mathrm{noise+anisotropic}}, \qquad
\mathcal{M}_{\mathrm{noise+fixed}\text{-}L}.
\end{equation}
The key quantities to report are log-evidence differences,
\begin{equation}
\Delta \log Z_{ab}=\log Z_a-\log Z_b.
\end{equation}

\begin{table}[H]
\centering
\caption{Placeholder evidence table. Replace with the final BLIP log-evidence values and Bayes-factor comparisons.}
\label{tab:evidence}
\begin{tabular}{lll}
\toprule
Model & $\log Z$ & Comment \\
\midrule
Noise only & placeholder & placeholder \\
Noise + isotropic SGWB & placeholder & placeholder \\
Noise + anisotropic SGWB & placeholder & placeholder \\
Noise + fixed-$L$ SGWB & placeholder & placeholder \\
\bottomrule
\end{tabular}
\end{table}

Placeholder text: discuss whether the anisotropic model is preferred over the isotropic alternative, whether the evidence gain justifies the extra complexity of the full spherical-harmonic model, and whether the single-multipole mode is useful as a screening analysis before a more expensive full reconstruction.

\section{Conclusion}
BLIP provides a coherent Bayesian framework for anisotropy inference in LISA SGWB analyses. The method combines physically motivated LISA noise spectra, time-dependent orbital responses, configurable spectral models, and covariance-based likelihood evaluation in either a full spherical-harmonic or reduced fixed-multipole description. The most important implementation feature is that anisotropy enters through the detector covariance itself, so the inference naturally uses both the auto-channel power and the complex cross-channel structure induced by the moving detector.

For a final paper, the placeholders in Section~3 should be replaced with production-run posterior summaries, model-fit diagnostics, and evidence comparisons. If the aim is a first detection claim for anisotropy at a given angular scale, the fixed-$L$ total-power mode offers a concise result. If the aim is directional reconstruction, the full spherical-harmonic mode provides the appropriate map-level inference.

\appendix
\section{Code-to-Method Mapping}
For reproducibility, the main implementation pieces described above are currently located in the following BLIP modules:
\begin{itemize}
\item \texttt{blip/run\_blip}: configuration parsing, data generation, sampler dispatch, and output writing;
\item \texttt{blip/src/makeLISAdata.py}: time-series handling, SGWB injection, and Fourier-domain data products;
\item \texttt{blip/src/instrNoise.py}: instrumental noise spectra and noise realizations;
\item \texttt{blip/src/geometry.py} and \texttt{blip/src/sph\_geometry.py}: isotropic and anisotropic detector responses;
\item \texttt{blip/src/clebschGordan.py}: conversion between sampled $b_{\ell m}$ coefficients and power-map harmonics;
\item \texttt{blip/src/models.py}: prior transforms, covariance assembly, and the unified likelihood;
\item \texttt{blip/src/dynesty\_engine.py}, \texttt{blip/src/nessai\_engine.py}, and \texttt{blip/src/emcee\_engine.py}: posterior samplers.
\end{itemize}
If this manuscript is used as the basis for a methods section in a publication, version-controlling the exact BLIP commit hash alongside the final output directory is strongly recommended.

\end{document}
